% Part: incompleteness
% Chapter: representability-in-q
% Section: c-representable

\documentclass[../../../include/open-logic-section]{subfiles}

\begin{document}

\olfileid{inc}{req}{cre}
\olsection{Fungsi ing $C$ Bisa Direpresentasikaké ing $\Th{Q}$}

Kita kudu nuduhaké yèn saben fungsi ing $C$ bisa direpresentasikaké
ing $\Th{Q}$. Pungkasané, kita kudu nuduhaké carané kanggo saben
fungsi aritas-$k$ $f(x_0,\dots,x_{k-1})$ ing $C$ milih formula
$!A_f(x_0,\dots,x_{k-1},y)$ kang merepresentasikaké fungsi mau.

\begin{lem}
Kanggo wilangan asli $n$ lan $m$, yèn $n \neq m$, mula
$Q \Proves \num n \neq \num m$.
\end{lem}

\begin{proof}
Gunakna induksi tumrap $n$ kanggo nuduhaké yèn kanggo saben $m$,
yèn $n \neq m$, mula $Q \Proves \eq/[\num n][\num m]$.

Ing kahanan dhasar, $n = 0$. Yèn $m$ ora padha karo $0$,
mula $m = k + 1$ kanggo sawatara wilangan asli $k$.
Kita nduwèni aksioma $\lforall[x][\eq/[0][x']]$.
Kanthi aksioma kuantor, kita ganti $x$ nganggo $\num k$
lan olèh $\eq/[0][\num k']$. Nanging $\num k'$ iku
$\num m$.

Ing langkah induksi, anggepen pratelan iki bener kanggo $n$,
banjur delengen $n+1$. Anggep $m$ wilangan asli sembarang.
Ana rong kemungkinan: $m = 0$ utawa kanggo sawatara $k$
kita nduwèni $m = k+1$. Kahanan kapisan dirampungaké kaya ing ndhuwur.
Ing kahanan kapindho, anggepen $n+1 \neq k+1$.
Mula $n \neq k$. Miturut hipotesis induksi kanggo $n$,
kita olèh $\Th{Q} \Proves \eq/[\num n][\num k]$.
Ana aksioma $\lforall[x][\lforall[y][\eq[x'][y'] \lif \eq[x][y]]]$.
Kanthi aksioma kuantor, kita olèh
$\eq[\num n'][\num k'] \lif \eq[\num n][\num k]$.
Miturut logika proposisional, ing $\Th{Q}$ kita bisa ngasilaké
$\eq/[\num n][\num k] \lif \eq/[\num n'][\num k']$.
Kanthi modus ponens, kita olèh $\eq/[\num n'][\num k']$,
kaya kang dikarepaké, amarga $\num k'$ iku $\num m$.
\end{proof}

Lema iki pratelané ora akèh: intiné, $\Th{Q}$ bisa mbuktèkaké
yèn numeral kang béda nuduhaké obyek kang béda. Umpamané,
$\Th{Q}$ mbuktèkaké $0'' \neq 0'''$. Nanging bukti umumé
butuh pangati-ati. Élinga uga yèn induksi kang kita gunakaké
ana \emph{ing njaba} $\Th{Q}$.

Kita bakal bisa merepresentasikaké nol, penerus,
panjumlahan, perkalian, fungsi karakteristik kanggo
kesamaan, lan proyeksi. Ing saben kahanan, formula kang
merepresentasikaké fungsi iku gampang ditemtokaké;
umpamané, nol direpresentasikaké déning formula
\[
y = 0,
\]
penerus direpresentasikaké déning formula
\[
x_0' = y,
\]
lan panjumlahan direpresentasikaké déning formula
\[
x_0 + x_1 = y.
\]
Pakaryan kang kudu ditindakake yaiku nuduhaké yèn $\Th{Q}$
bisa mbuktèkaké pratelan kang gegayutan. Umpamané, yèn
panjumlahan direpresentasikaké déning formula ing ndhuwur,
kita kudu nuduhaké yèn kanggo saben pasangan wilangan asli
$m$ lan $n$, $\Th{Q}$ mbuktèkaké
\[
\eq[\num n + \num m][\num {n+m}]
\]
lan
\[
\lforall[y][(\eq[(\num n + \num m)][y] \lif \eq[y][\num {n+m}])].
\]

Kepiyé komposisi? Anggepen $h$ ditetepaké déning
\[
h(x_0,\dots,x_{l-1}) = f(g_0(x_0,\dots,x_{l-1}), \dots,
g_{k-1}(x_0,\dots,x_{l-1})).
\]
ing kéné kita wis nemu formula $!A_f, !A_{g_0}, \dots,
!A_{g_{k-1}}$ kang merepresentasikaké fungsi
$f,g_0,\dots,g_{k-1}$ miturut urutané. Banjur kita bisa
netepaké formula $!A_h$ kang merepresentasikaké $h$,
yaiku kanthi netepaké $!A_h(x_0,\dots,x_{l-1},y)$ dadi
\begin{multline*}
  \lexists[z_0\dots][\lexists[z_{k-1}][(!A_{g_0}(x_0,\dots,x_{l-1},z_0) \land
  \dots \land !A_{g_{k-1}}(x_0,\dots,x_{l-1},z_{k-1}) \land\\
  !A_f(z_0,\dots,z_{k-1},y))]].
\end{multline*}

Pungkasané, delengen telusuran tanpa wates. Anggepen
$g(x,\vec z)$ reguler lan bisa direpresentasikaké ing $\Th{Q}$,
umpamané déning formula $!A_g(x,\vec z,y)$.
Anggep $f$ ditetepaké déning $f(\vec z) = \mu x \; g(x,\vec z)$.
Kita arep nemu formula $!A_f(\vec z,y)$ kang merepresentasikaké
$f$. Pilihan kang lumrah yaiku:
\[
!A_f(\vec z,y) \ident !A_g(y,\vec z,0) \land \lforall[w][(w < y
\lif \lnot !A_g(w,\vec z,0))].
\]

Bisa dituduhaké yèn iki bener kanthi sawetara lema tambahan,
umpamané:

\begin{lem}
  Kanggo saben variabel $x$ lan wilangan asli $n$,
  $\Th{Q}$ mbuktèkaké $\eq[(x' + \num n)][(x + \num n)']$.
\end{lem}

Élinga manèh yèn pratelan iki luwih lemah tinimbang
pratelan yèn $\Th{Q}$ mbuktèkaké
$\lforall[x][\lforall[y][(x' + y) = (x +y)']]$
(kang pancèn salah).

\begin{proof}
Buktiné, kaya lumrahé, nganggo induksi tumrap $n$.
Ing kahanan dhasar, $n = 0$, kita kudu nuduhaké yèn
$\Th{Q}$ mbuktèkaké $\eq[(x'+0)][(x+0)']$. Nanging kita olèh:
\begin{eqnarray*}
& & \eq[(x' + 0)][x'] \quad \text{saka aksioma 4} \\
& & \eq[(x + 0)][x] \quad \text{saka aksioma 4} \\
& & \eq[(x+0)'][x'] \quad \text{miturut larik 2} \\
& & \eq[(x' + 0)][(x + 0)'] \quad \text{larik 1 lan 3}
\end{eqnarray*}
Ing langkah induksi, anggepen kita wis nurunké $\eq[(x' +
  \num n)][(x + \num n)']$ ing $\Th{Q}$.
Amarga $\num{n+1}$ iku $\num n'$, kita kudu nuduhaké yèn
$\Th{Q}$ mbuktèkaké $\eq[(x' + \num n')][(x + \num n')']$.
Rerangkening kesamaan ing ngisor iki bisa diturunké ing $\Th{Q}$:
\begin{eqnarray*}
(x' + \num n') & \eq & (x' + \num n)' \quad \text{aksioma 5} \\
& \eq & (x + \num n')' \quad \text{saka hipotesis induksi}
\end{eqnarray*}
\end{proof}

\begin{lem}
\begin{enumerate}
\item $\Th{Q}$ mbuktèkaké $\lnot (x < \num 0)$.
\item Kanggo saben wilangan asli $n$, $\Th{Q}$ mbuktèkaké
\[
x < \num {n+1} \lif (\eq[x][\num 0] \lor \dots \lor \eq[x][\num n]).
\]
\end{enumerate}
\end{lem}

\begin{proof}
Kita rampungaké pérangan 1 lan sawatara saka pérangan 2 kanthi
informal (yaiku mung mènèhi pituduh carané mbangun
!!{derivation} formal).

Kanggo pérangan 1, miturut definisi $<$, kita kudu mbuktèkaké
$\lnot \lexists[y][\eq[(y'+x)][0]]$ ing $\Th{Q}$.
Miturut aksioma lan aturan logika tataran kapisan, iki padha
karo $\lforall[y][\eq/[(y' + x)][0]]$.
Gagasané mangkéné: anggepen $\eq[(y'+x)][0]$.
Yèn $x$ iku 0, kita olèh $\eq[(y' + 0)][0]$.
Nanging miturut aksioma 4 saka $\Th{Q}$, kita olèh
$\eq[(y' + 0)][y']$, lan miturut aksioma 2 kita olèh
$\eq/[y'][0]$, sawijining kontradiksi.
Mula $\lforall[y][\eq/[(y' +x)][0]]$.
Yèn $x$ dudu $0$, miturut aksioma 3 ana $z$ saéngga
$\eq[x][z']$. Banjur kita olèh $\eq[(y' + z')][0]$.
Miturut aksioma~5, kita olèh $\eq[(y' + z)'][0]$,
kang manèh mbantah aksioma~2.

Kanggo pérangan 2, gunakna induksi tumrap $n$.
Delengen kahanan dhasar nalika $n =0$.
Ing kahanan iki kita kudu nuduhaké
$x < \num 1 \lif \eq[x][\num 0]$.
Anggepen $x < \num 1$. Miturut aksioma kang netepaké $<$,
kita olèh $\lexists[y][\eq[(y'+x)][0']]$.
Anggepen $y$ nduwèni sipat iki; mula
$\eq[y'+x][0']$.

Kita kudu nuduhaké $\eq[x][0]$. Miturut aksioma 3,
yèn $x$ dudu 0, mula padha karo $z'$ kanggo sawatara $z$.
Banjur kita olèh $\eq[(y' + z')][0']$.
Miturut aksioma~5 saka $\Th{Q}$, kita olèh
$\eq[(y'+z)'][0']$. Miturut aksioma 1, kita olèh
$\eq[(y' + z)][0]$. Nanging miturut definisi,
iki ateges $z < 0$, kang mbantah pérangan~1.
\end{proof}

\begin{explain}
Kita wis nuduhaké yèn himpunan fungsi kang bisa diitung bisa
ditandhani minangka himpunan fungsi kang bisa direpresentasikaké
ing $\Th{Q}$. Sejatiné, buktiné luwih umum.
Saka definisi representabilitas, gampang dideleng yèn
saben téori kang ngluwihi $\Th{Q}$ (utawa kang bisa
nginterpretasikaké $\Th{Q}$) bisa merepresentasikaké
fungsi-fungsi kang bisa diitung. Kosok baliné, ing saben
sistem !!{derivation} kang konsep buktiné bisa diitung,
saben fungsi kang bisa direpresentasikaké uga bisa diitung.
Mula, umpamané, himpunan fungsi kang bisa diitung bisa
ditandhani minangka himpunan fungsi kang direpresentasikaké
ing aritmetika Peano, utawa malah ing téori himpunan Zermelo
Fraenkel. Kaya kang dicathet G\"odel, iki rada nggumunaké.
Bakal kita deleng yèn pitakonan bab apa kang bisa dibuktèkaké
gumantung banget marang téori kang dipilih; kanthi kasar,
saya kuwat aksiomané, saya akèh kang bisa dibuktèkaké.
Nanging ing akèh téori aksiomatik, fungsi kang bisa
direpresentasikaké persis fungsi kang bisa diitung.
\end{explain}

\end{document}
